\section{Results}
\label{sec:results}

\subsection{Test bed: Minnesota, 2020 PL 94-171}

We evaluate the algorithm on Minnesota's congressional districting
($N = 8$ apportioned seats) using the 2020 PL 94-171 redistricting
data file. The input is the TIGER/Line 2020 VTD (voting tabulation
district) shapefile for Minnesota (4{,}110 atomic units, total
population $P = 5{,}706{,}494$, total land area $A = 225{,}187$~km$^2$),
joined to the Census Data API's \verb|P1_001N| total-population
field, projected to EPSG:5070 (CONUS Albers, equal area) at load
time. Per-district targets are
$P^{*} = P/8 = 713{,}312$ and $A^{*} = A/8 = 28{,}148$~km$^2$. All
runs are deterministic and reproducible from the script
\verb|scripts/prep_mn_units.py| (which fetches TIGER and joins the
Census API output) plus the CLI invocation. The data pipeline is
documented end-to-end in \verb|docs/mn-poc-walkthrough.md|.

We compare three plans on the same 4{,}110-VTD input:

\begin{itemize}
  \item \textbf{Pure radial.} The default algorithm
        (\S\ref{sec:methods-seeds}--\S\ref{sec:methods-repair}). No
        tuning knobs.
  \item \textbf{Radial $+$ tighten-pop ($\tau = 0.005$).} The default
        algorithm followed by the optional L$^1$ population-tightening
        pass (\S\ref{sec:methods-tighten}) targeting per-district
        deviation within 0.5\,\% of $P^{*}$.
  \item \textbf{Enacted Minnesota 119th-Congress plan.} The
        court-drawn, legislatively enacted U.S.~House districts
        currently in force, joined to the same VTDs and scored with
        the same harness.
\end{itemize}

\subsection{Headline scores}

Table~\ref{tab:mn-results} reports the DualBalance Score and supporting
metrics for the three plans.

\begin{table}[h]
\centering
\caption{DualBalance metrics on the Minnesota PoC (4{,}110 VTDs,
8~districts, 2020 PL 94-171 population). $\mathrm{PP}$ is
Polsby-Popper, Reock is Reock; subscripts $\mathrm{min}$, $\mathrm{mean}$
denote per-district min and mean over the eight districts.
$\mathrm{pop\_dev}$, $\mathrm{area\_dev}$ are the relative deviations
defined in \S\ref{sec:methods}. Best result per row in bold.}
\label{tab:mn-results}
\begin{tabular}{lrrr}
\toprule
Metric & Pure radial & Radial $+$ tighten-pop & Enacted (119th) \\
\midrule
$\mathrm{DBS}$ & 0.6472 & \textbf{0.6574} & 0.6390 \\
$\overline{\mathrm{pop\_dev}}$ & 5.08\,\% & \textbf{0.08\,\%} & 0.42\,\% \\
$\mathrm{pop\_dev}_{\max}$ & 11.24\,\% & \textbf{0.21\,\%} & 1.32\,\% \\
$\overline{\mathrm{area\_dev}}$ & \textbf{103.9\,\%} & 104.2\,\% & 112.6\,\% \\
$\mathrm{area\_dev}_{\max}$ & 271.0\,\% & 275.1\,\% & \textbf{241.0\,\%} \\
$\mathrm{PP}_{\mathrm{mean}}$ & 0.200 & 0.162 & \textbf{0.320} \\
$\mathrm{PP}_{\min}$ & 0.094 & 0.047 & \textbf{0.178} \\
$\mathrm{Reock}_{\mathrm{mean}}$ & 0.361 & 0.342 & \textbf{0.419} \\
\bottomrule
\end{tabular}
\end{table}

Three findings:

\paragraph{Pure radial beats the enacted plan on the DualBalance Score.}
The pure-radial pipeline scores 0.6472 against the enacted plan's
0.6390, a margin of $+1.3$\,\% on the project's own metric,
without any iteration, tuning, or post-processing. The advantage comes
entirely from area balance ($\overline{\mathrm{area\_dev}}$ 103.9\,\%
vs 112.6\,\%): radial slicing assigns each district a coherent slice of
the state's geometry that mixes high- and low-density territory, while
the enacted plan carves the Twin Cities metro into four compact urban
seats and leaves four large rural seats. The DualBalance Score
penalizes the latter pattern more heavily because it weights area
deviation equally with population deviation.

\paragraph{The optional tighten-pop pass closes the legal gap and
improves the DualBalance Score.} The pure-radial plan's
$\mathrm{pop\_dev}_{\max}$ of 11.24\,\% is well above the
\textasciitilde{}0.5\,\% deviation typical of enacted U.S.~House
plans~\cite{reynolds1964}. Running the optional L$^1$ pop-tightening
pass (\S\ref{sec:methods-tighten}) executed 81 boundary-unit swaps in
\textasciitilde{}18\,seconds of compute, driving
$\mathrm{pop\_dev}_{\max}$ from 11.24\,\% to 0.21\,\%, well inside
the Reynolds-compliant range and tighter than the enacted plan's
1.32\,\%. $\overline{\mathrm{area\_dev}}$ rose by only 0.3 percentage
points; the DualBalance Score \emph{improved} from 0.6472 to 0.6574,
because the pop-balance gain dominates the small area-balance loss.
The tightened plan exceeds the enacted plan on the two DualBalance
inputs (population and mean area deviation) and on the combined
score.

\paragraph{Compactness is the price.} The enacted plan beats DualBalance
on every compactness measure: $\mathrm{PP}_{\min}$ 0.178 vs 0.094 in
the radial plan and 0.047 after pop-tightening, $\mathrm{PP}_{\mathrm{mean}}$
0.320 vs 0.200 / 0.162, $\mathrm{Reock}_{\mathrm{mean}}$ 0.419 vs
0.361 / 0.342. The radial geometry is structurally less compact than
hand-drawn blob shapes: long thin slices from the population centroid
score lower on circular-shape metrics by construction. Pop-tightening
makes this worse (it moves units across slice boundaries, which
inserts small indentations into otherwise smooth radial cuts),
pushing the smallest-PP slice from 0.094 (already at the
``0.10 raises eyebrows'' informal threshold) down to 0.047. Whether
the trade is acceptable is a normative question we return to in
\S\ref{sec:discussion-compactness}; on the algorithm's own metric,
the trade is favorable.

\subsection{Map comparison}

Figure~\ref{fig:mn-comparison} displays the three plans side by side.

\begin{figure}[h]
\centering
\includegraphics[width=\textwidth]{mn_radial_with_tighten.png}
\caption{Minnesota congressional districts under three plans.
\textbf{Left:} the enacted 119th-Congress plan, with four compact
Twin Cities seats and four large rural seats.
\textbf{Center:} pure-radial DualBalance (score $0.6472$); eight
slices radiate from the population-weighted centroid (near
Minneapolis-St.\ Paul), each spanning both dense and sparse
territory. \textbf{Right:} radial $+$ \texttt{--tighten-pop 0.005}
(score $0.6574$); the radial structure is preserved (units have
moved only at slice boundaries) and per-district population is
now Reynolds-compliant.}
\label{fig:mn-comparison}
\end{figure}

\subsection{Determinism check}

We re-ran \verb|dualbalance generate --config configs/mn_vtd.yaml|
ten times in succession, comparing each run's \verb|map.geojson| and
\verb|metrics.json| by byte hash. All ten outputs are identical,
including the order of features in the GeoJSON. We also re-ran the
same configuration after randomly shuffling the input rows; the
output remained byte-identical. The CLI integration test
\verb|test_generate_determinism_via_cli| pins this property in CI.

\subsection{Per-district breakdown}

Table~\ref{tab:mn-perdistrict} reports the per-district metrics for
the radial $+$ tighten-pop plan, indexed by seed angle (district 0
sits due east of the population-weighted centroid; subsequent
districts advance counter-clockwise in steps of $2\pi/8$ rad).
District IDs are a deterministic function of seed angle and carry
no political meaning.

\begin{table}[h]
\centering
\caption{Per-district breakdown, radial $+$ tighten-pop 0.5\,\%.
Areas are in km$^2$.}
\label{tab:mn-perdistrict}
\begin{tabular}{rrrrrrr}
\toprule
District & Population & Area & $\mathrm{pop\_dev}$ & $\mathrm{area\_dev}$ & $\mathrm{PP}$ & $\mathrm{Reock}$ \\
\midrule
0 & 712{,}841 &   1{,}010 & 0.07\,\% &  96.4\,\% & 0.228 & 0.454 \\
1 & 714{,}816 &   9{,}620 & 0.21\,\% &  65.8\,\% & 0.047 & 0.147 \\
2 & 712{,}145 & 105{,}594 & 0.16\,\% & 275.1\,\% & 0.152 & 0.369 \\
3 & 713{,}263 &  50{,}054 & 0.01\,\% &  77.8\,\% & 0.146 & 0.268 \\
4 & 714{,}168 &  46{,}060 & 0.12\,\% &  63.6\,\% & 0.150 & 0.267 \\
5 & 713{,}368 &  10{,}882 & 0.01\,\% &  61.3\,\% & 0.137 & 0.422 \\
6 & 713{,}013 &   1{,}481 & 0.04\,\% &  94.7\,\% & 0.117 & 0.301 \\
7 & 712{,}880 &      477 & 0.06\,\% &  98.3\,\% & 0.318 & 0.510 \\
\bottomrule
\end{tabular}
\end{table}

District 2 inherits Minnesota's northern panhandle and is the largest
by area ($\sim$3.8$\times A^{*}$), while District 7 contains only the
Twin Cities urban core ($\sim$0.02$\times A^{*}$). All eight districts
are within 0.21\,\% of $P^{*}$. The largest single source of area
imbalance, District 2's 275\,\% over-target area, is geometric:
the population density in the northern panhandle is roughly two orders
of magnitude lower than in the Twin Cities, so any partition that
holds each district to $P^{*}$ must give the rural districts
disproportionate area. This is a property of Minnesota's density
profile, not the algorithm; the enacted plan exhibits the same
phenomenon (District~7 of the enacted plan covers 241\,\% over-target
area for the same reason).

\subsection{Descriptive diagnostics: partisan, race, county}
\label{sec:results-diagnostics}

We report a second tier of numbers (partisan, racial, and
county-split metrics) not because the algorithm uses them, but
because readers will reasonably ask how a procedurally-neutral plan
\emph{looks} on the diagnostics conventionally applied to enacted
maps. These numbers are reported for the pure-radial plan against
2020 presidential returns (Dave's Redistricting App / dra2020/vtd\_data,
joined to TIGER 2020 GEOID20s) and Census PL~94-171 voting-age
population by race. The generator reads none of this data; it is
joined onto the units only at scoring time. We emphasize, in line
with \S\ref{sec:discussion-forensic}, that these are descriptive
statistics of the resulting partition, not evaluative scores of it.

\paragraph{Partisan profile.} Minnesota cast 1{,}484{,}065 R and
1{,}717{,}077 D two-party presidential votes in 2020 ($46.4\,\%$
statewide R share, $53.6\,\%$ D). On the pure-radial plan the eight
districts split 4--4 between R-winning and D-winning, against a
seats-proportional expectation of $3.7$ R seats. The two D-leaning
slices contain the Minneapolis-St.~Paul metro and run at 64--75\,\% D;
the two R-leaning slices reaching into the western and southern parts
of the state run at 55--65\,\% R; the four remaining slices are within
$\pm 5$ points of even. The efficiency gap is $+6.6\,\%$ (positive =
R-favorable) and the mean-median R difference is $+3.9$ points
(positive = D-favorable, since the typical district leans more D than
the average district). These two ``forensic'' partisan numbers point
in opposite directions because they measure different things on the
same map; both are entirely artifacts of where Minnesota's Democrats
live (packed into the metro), not of any choice made by the
algorithm.

\paragraph{Race / VAP profile.} Minnesota's voting-age population is
80.0\,\% non-Hispanic white, 5.9\,\% Black, 5.0\,\% Hispanic, 4.9\,\%
Asian, and 1.1\,\% American Indian / Alaska Native. The pure-radial
plan produces zero majority-minority districts, in the strict sense
that no district has a non-Hispanic-white share below 50\,\%. The
district containing the Minneapolis-St.~Paul urban core (District~7)
has the lowest non-Hispanic-white share at 65.3\,\%, the highest
Black VAP share at 15.7\,\%, and the highest Hispanic VAP share at
6.7\,\%; the remaining seven districts run 73--90\,\% non-Hispanic
white. The enacted 119\textsuperscript{th}-Congress Minnesota plan
also has no majority-minority districts. Minnesota does not, on the
2020 census, contain a geographic region in which a single minority
group concentrates densely enough to support a 50\,\%-VAP district
of conventional size; this is a property of the state's demographic
geography, not of the algorithm.

\paragraph{County splits.} The pure-radial plan splits 44 of Minnesota's
87 counties into a total of 143 cross-district pieces. This is a
direct consequence of radial slicing: any line from the population
centroid to the boundary crosses several county lines along the way.
The enacted plan splits fewer counties (it is drawn to follow them
more closely), at the cost of the larger area imbalance that
Table~\ref{tab:mn-results} reports. Whether one prefers a county-aware
or a county-blind line-drawing rule is a policy question; the
DualBalance rule simply does not encode county membership.

\paragraph{What to read into these numbers.} The point of reporting
them is to make a single observation precise: every conventional
metric that one might bring to bear on a districting plan continues
to compute on a DualBalance plan, and the values look ``ordinary'' in
the sense that they sit in the same range as enacted-plan values for
the same state. But the metrics' \emph{interpretation} differs by
construction. On an enacted plan, an efficiency gap of $+6.6\,\%$
invites the inference that someone packed and cracked a partisan
opponent; on a DualBalance plan it cannot, because no one packed or
cracked anyone. On an enacted plan, zero majority-minority districts
invites the inference that the line-drawers either declined to honor
\textit{Allen v.~Milligan}~\cite{milligan2023} or did honor it on a
geography that could not support a majority-minority seat; on a
DualBalance plan only the latter inference is available, because
the algorithm has no way of seeing race. The numbers are the same;
what they license is different. We return to this point in
\S\ref{sec:discussion-forensic}.
